\documentclass[10pt,a4paper]{amsart}
\usepackage{amsmath,amssymb,amsthm,graphicx}
\usepackage{verbatim}
\newenvironment{reference}{\comment}{\endcomment}
%\newenvironment{reference}{}{}
\newenvironment{slogan}{\comment}{\endcomment}
\newenvironment{history}{\comment}{\endcomment}

% For commutative diagrams we use Xy-pic
\usepackage[all]{xy}

% We use 2cell for 2-commutative diagrams.
\xyoption{2cell}
% Theorem environments.
%
\theoremstyle{plain}
\newtheorem{theorem}[subsection]{Theorem}
\newtheorem{proposition}[subsection]{Proposition}
\newtheorem{lemma}[subsection]{Lemma}

\theoremstyle{definition}
\newtheorem{definition}[subsection]{Definition}
\newtheorem{example}[subsection]{Example}
\newtheorem{exercise}[subsection]{Exercise}
\newtheorem{situation}[subsection]{Situation}

\theoremstyle{remark}
\newtheorem{remark}[subsection]{Remark}
\newtheorem{remarks}[subsection]{Remarks}

\numberwithin{equation}{subsection}

% Macros
%
\def\lim{\mathop{\mathrm{lim}}\nolimits}
\def\colim{\mathop{\mathrm{colim}}\nolimits}
\def\Spec{\mathop{\mathrm{Spec}}}
\def\Hom{\mathop{\mathrm{Hom}}\nolimits}
\def\Ext{\mathop{\mathrm{Ext}}\nolimits}
\def\SheafHom{\mathop{\mathcal{H}\!\mathit{om}}\nolimits}
\def\SheafExt{\mathop{\mathcal{E}\!\mathit{xt}}\nolimits}
\def\Sch{\mathit{Sch}}
\def\Mor{\mathop{\mathrm{Mor}}\nolimits}
\def\Ob{\mathop{\mathrm{Ob}}\nolimits}
\def\Sh{\mathop{\mathit{Sh}}\nolimits}
\def\NL{\mathop{N\!L}\nolimits}
\def\CH{\mathop{\mathrm{CH}}\nolimits}
\def\proetale{{pro\text{-}\acute{e}tale}}
\def\etale{{\acute{e}tale}}
\def\QCoh{\mathit{QCoh}}
\def\Ker{\mathop{\mathrm{Ker}}}
\def\Im{\mathop{\mathrm{Im}}}
\def\Coker{\mathop{\mathrm{Coker}}}
\def\Coim{\mathop{\mathrm{Coim}}}

% Boxtimes
%
\DeclareMathSymbol{\boxtimes}{\mathbin}{AMSa}{"02}

%
% Macros for moduli stacks/spaces
%
\def\QCohstack{\mathcal{QC}\!\mathit{oh}}
\def\Cohstack{\mathcal{C}\!\mathit{oh}}
\def\Spacesstack{\mathcal{S}\!\mathit{paces}}
\def\Quotfunctor{\mathrm{Quot}}
\def\Hilbfunctor{\mathrm{Hilb}}
\def\Curvesstack{\mathcal{C}\!\mathit{urves}}
\def\Polarizedstack{\mathcal{P}\!\mathit{olarized}}
\def\Complexesstack{\mathcal{C}\!\mathit{omplexes}}
% \Pic is the operator that assigns to X its picard group, usage \Pic(X)
% \Picardstack_{X/B} denotes the Picard stack of X over B
% \Picardfunctor_{X/B} denotes the Picard functor of X over B
\def\Pic{\mathop{\mathrm{Pic}}\nolimits}
\def\Picardstack{\mathcal{P}\!\mathit{ic}}
\def\Picardfunctor{\mathrm{Pic}}
\def\Deformationcategory{\mathcal{D}\!\mathit{ef}}
\usepackage{cleveref}
\crefname{originalreference}{original reference}{original references}
\Crefname{originalreference}{Original reference}{Original references}
\makeatletter
\newlabel{dpa-definition-divided-powers}{{07GL}{source}{}{sourcepointer-1}{}}
\newlabel{dpa-definition-divided-powers@cref}{{[originalreference][0][]07GL}{[1][0][]source}{}{}{}}
\makeatother
\crefname{originalreference}{original reference}{original references}
\Crefname{originalreference}{Original reference}{Original references}
\makeatletter
\newlabel{dpa-lemma-silly}{{07GM}{source}{}{sourcepointer-2}{}}
\newlabel{dpa-lemma-silly@cref}{{[originalreference][0][]07GM}{[1][0][]source}{}{}{}}
\makeatother
\title{Divided power structure on the second-order thickening A \textbackslash{}\textbackslash{}oplus M \textbackslash{}\textbackslash{}oplus N}
\author{The Stacks Project Authors}
\date{stacks-project:crystalline.tex:local-snapshot-20261012; GNU FDL 1.2 or later}
\begin{document}
\maketitle
\noindent\textit{Editorial note.} All delivered mathematical text is the author's own, copied verbatim from the pinned Stacks Project chapter file crystalline.tex; no statement or proof step was rewritten, reordered or removed.\par\medskip
\noindent\textit{Editorial note.} Span recomputation by label pairing (LANE\_BRIEF 10.2): in this chapter the author's proof environment follows the statement environment as a standalone block; the spec covers the complete environment pair (statement span, proof span) as recorded in the pool candidate.\par\medskip
\ifdefined\hypertarget\hypertarget{sourcepointer-1}{}\fi
\noindent\textit{Source reference.} Original reference 07GL: sources/stacks-chow/tags.txt, source lines 9521--9521. Cross-chapter standard prerequisite cited by the author inside the counted span; it is used as a pointer and its own proof is not part of this item.\par\medskip
\ifdefined\hypertarget\hypertarget{sourcepointer-2}{}\fi
\noindent\textit{Source reference.} Original reference 07GM: sources/stacks-chow/tags.txt, source lines 9522--9522. Cross-chapter standard prerequisite cited by the author inside the counted span; it is used as a pointer and its own proof is not part of this item.\par\medskip
\begin{lemma}
\label{lemma-divided-power-second-order-thickening}
Let $(A, I, \gamma)$ be a divided power ring. Let $M$, $N$ be $A$-modules.
Let $q : M \times M \to N$ be an $A$-bilinear map.
Let $B = A \oplus M \oplus N$ as an $A$-algebra with multiplication
$$
(x, z, w)\cdot (x', z', w') = (xx', xz' + x'z, xw' + x'w + q(z, z') + q(z', z))
$$
and set $J = I \oplus M \oplus N$. Set
$$
\delta_n(x, z, w) = (\gamma_n(x), \gamma_{n - 1}(x)z,
\gamma_{n - 1}(x)w + \gamma_{n - 2}(x)q(z, z))
$$
for $(x, z, w) \in J$.
Then $\delta$ is a divided power structure and
$A \to B$ is a homomorphism of divided power rings from
$(A, I, \gamma)$ to $(B, J, \delta)$.
\end{lemma}

\begin{proof}
Suppose we want to prove that property (4) of
Divided Power Algebra, Definition \ref{dpa-definition-divided-powers}
is satisfied. Pick $(x, z, w)$ and $(x', z', w')$ in $J$.
Pick a map
$$
A_0 = \mathbf{Z}\langle s, s'\rangle \longrightarrow A,\quad
s \longmapsto x,
s' \longmapsto x'
$$
which is possible by the universal property of divided power
polynomial rings. Set $M_0 = A_0 \oplus A_0$ and
$N_0 = A_0 \oplus A_0 \oplus M_0 \otimes_{A_0} M_0$.
Let $q_0 : M_0 \times M_0 \to N_0$ be the obvious map.
Define $M_0 \to M$ as the $A_0$-linear map which sends
the basis vectors of $M_0$ to $z$ and $z'$. Define $N_0 \to N$
as the $A_0$ linear map which sends the first two basis vectors
of $N_0$ to $w$ and $w'$ and uses
$M_0 \otimes_{A_0} M_0 \to M \otimes_A M \xrightarrow{q} N$
on the last summand. Then we see that it suffices to prove the
identity (4) for the situation $(A_0, M_0, N_0, q_0)$.
Similarly for the other identities. This reduces us to the case
of a $\mathbf{Z}$-torsion free ring $A$ and $A$-torsion free modules.
In this case all we have to do is show that
$$
n! \delta_n(x, z, w) = (x, z, w)^n
$$
in the ring $A$, see Divided Power Algebra, Lemma \ref{dpa-lemma-silly}.
To see this note that
$$
(x, z, w)^2 = (x^2, 2xz, 2xw + 2q(z, z))
$$
and by induction
$$
(x, z, w)^n = (x^n, nx^{n - 1}z, nx^{n - 1}w + n(n - 1)x^{n - 2}q(z, z))
$$
On the other hand,
$$
n! \delta_n(x, z, w) = (n!\gamma_n(x), n!\gamma_{n - 1}(x)z,
n!\gamma_{n - 1}(x)w + n!\gamma_{n - 2}(x) q(z, z))
$$
which matches. This finishes the proof.
\end{proof}

\end{document}
